\documentclass[11pt,a4paper]{article}
\usepackage[etiket=en]{gtu-el-yazisi}
\begin{document}
\notbilgi{Elektronik II}{Feedback: General Structure, Voltage and Current Amplifiers}{Ders11.pdf}
\textit{Not: Orijinal dosyadaki basılı sayfalar (s. 2--5) dijitalleştirilmedi; yalnızca el yazısı sayfalar aktarıldı.}

\sayfa{1}
\section{Feedback}
\noindent
\begin{minipage}[c]{0.36\linewidth}
  \centering
  \begin{circuitikz}[line width=0.7pt, scale=0.85, transform shape]
    \draw (0,0) node[nmos] (M) {};
    \draw (M.G) to[short, -o] ++(-0.8,0) node[left] {$v_i$};
    \draw (M.D) to[R=$R_D$] ++(0,2.2) node[vcc] {$V_{DD}$};
    \draw (M.D) to[short, -o] ++(1.2,0) node[right] {$v_o$};
    \draw (M.S) to[R=$R_S$] ++(0,-2.2) node[ground] {};
  \end{circuitikz}
\end{minipage}\hfill
\begin{minipage}[c]{0.58\linewidth}
  Common source amplifier
  \[ \frac{v_o}{v_i} \approx \frac{R_D}{R_S} \]
\end{minipage}

\section{The general feedback structure}
\begin{center}
\begin{tikzpicture}[>=Stealth, line width=0.7pt, font=\small, box/.style={draw, minimum height=10mm}]
  \node[box, minimum width=22mm] (src) at (0,0) {Source};
  \node[draw, circle, inner sep=1pt] (sum) at (4,0) {$\Sigma$};
  \node[box, minimum width=18mm] (A) at (7,0) {$\overrightarrow{A}$};
  \node[box, minimum width=22mm] (load) at (12,0) {Load};
  \node[box, minimum width=18mm] (B) at (7,-1.8) {$\overleftarrow{\beta}$};
  \draw[->] (src) -- node[above] {$x_s$} (sum);
  \draw[->] (sum) -- node[above] {$x_i$} (A);
  \draw[->] (A) -- node[above, pos=0.8] {$x_o$} (load);
  \draw[->] (9.5,0) |- (B.east);
  \draw[->] (B.west) -| node[below, pos=0.2] {$x_f$} (sum);
  \node[anchor=south east, inner sep=0.5pt] at (sum.north west) {$+$};
  \node[anchor=north east, inner sep=0.5pt] at (sum.south west) {$-$};
\end{tikzpicture}
\end{center}

\noindent
\begin{minipage}[t]{0.47\linewidth}
  $A$ = open loop gain\\[1mm]
  $\beta$ = feedback factor\\[1mm]
  $x_i$ = Input signal\\[1mm]
  $x_s$ = signal source\\[1mm]
  $x_o$ = output signal\\[1mm]
  $x_f$ = sampled output by fbn (feedback network)
\end{minipage}\hfill
\begin{minipage}[t]{0.47\linewidth}
  $\left(\dfrac{x_o}{x_s}\right)_{CL} = \,?$\\[2mm]
  $\displaystyle\begin{aligned}
    x_o &= A\,x_i \\
    x_f &= \beta\cdot x_o \\
    x_i &= x_s - x_f \\
    x_i &= x_s - \beta\, x_o \\
    x_i &= x_s - A\beta\cdot x_i \\
    x_s &= x_i\,(1+A\beta)
  \end{aligned}$
\end{minipage}

\[ \left(\frac{x_o}{x_s}\right) = \frac{A\,x_i}{x_i\,(1+A\beta)} = \frac{A}{1+A\beta} \]

\noindent
\begin{minipage}[c]{0.5\linewidth}
  \centering
  \begin{circuitikz}[line width=0.7pt, scale=0.9, transform shape]
    \draw (0,0) node[op amp] (OA) {$A$};
    \coordinate (tap) at ($(OA.out)+(0.8,0)$);
    \draw (OA.-) -- ++(-0.6,0) coordinate (n1);
    \draw (n1) to[R=$R_1$] ++(-2,0) -- ++(0,-0.4) node[ground] {};
    \draw (n1) -- ++(0,1.5) coordinate (n2) to[R=$R_2$] (tap |- n2) -- (tap);
    \draw (OA.out) -- (tap) node[circ] {} to[short, -o] ++(1.2,0) node[right] {$v_o$};
    \draw (OA.+) -- ++(-0.6,0) to[sV, l=$v_i$] ++(0,-1.6) node[ground] {};
  \end{circuitikz}
\end{minipage}\hfill
\begin{minipage}[c]{0.46\linewidth}
  \[ \frac{v_o}{v_i} = \left(1+\frac{R_2}{R_1}\right) \textcolor{turuncukalem}{\;=\; \frac{1}{\beta}} \]
\end{minipage}

\sayfa{6}
\section{1- Voltage Amplifiers (Shunt (o) -- Series (I))}
\sekilnotu[6]{Block diagram of the shunt-series voltage amplifier: source $V_s$ with series $R_s$ feeds the ``Basic Amplifier'' block; the amplifier output is connected in parallel with the load $R_L$ and the output terminals ($+V_o-$); the ``Feedback'' block has its input ($+V_o-$) connected across the output (shunt), and its output voltage $V_f$ ($+$ top, $-$ bottom) is in series with the input loop (between $V_s$/$R_s$ and the amplifier input). Orange arrows mark $R_{if}$ (looking into the source loop), $R_{in}$ (looking into the amplifier input), $R_{out}$ (looking into the amplifier output, at $R_L$) and $R_{of}$ (looking into the output terminals).}

\section{Two port network equivalent}
\sekilnotu[6]{Two-port equivalent: source $V_s$ with $R_s$ (top) and $R_{11}$ (bottom wire; arrow labelled $R_i$ looking into the input) in front of the ``Basic Amplifier'' box; at the output $R_{22}$ and $R_L$ are in parallel across the output terminals ($+$ top, $-$ bottom, $V_o$; arrow labelled $R_o$ looking in from the output). A curved arrow pointing at the output carries the note ``feed forward gain $A = v_o/v_s$''.}

\sekilnotu[6]{Two sketches of the red ``Feedback'' box. Left: input terminals ($+$, $-$) open, output terminals shorted together, arrow labelled $R_{11}$ at the input side. Right: input terminals open, arrow labelled $R_{22}$ at the output side.}

To find $\beta$

\sekilnotu[6]{``Feedback'' box with current $I_1$ entering the input and $V_f$ ($+$, $-$) across the input terminals; a voltage source $V_o$ is connected across the output terminals. A red arrow pointing left (from output towards input) is drawn above the box.}

\formulkutu{\beta = \frac{V_f}{V_o}\Big|_{I_1=0}}

\sayfa{7}
\section{Complete model}
\sekilnotu[7]{Complete model circuit: source $V_s$ with $R_s$ feeds the input resistance $R_i$ (voltage $v_i$, $+$ top, $-$ bottom); the output side has the dependent source $A v_i$ in series with $R_o$ and the output terminals ($+v_o-$). The feedback part is a dependent voltage source $\beta V_o$ ($+$ top) in series with the input loop (between source terminal $s'$ and the bottom of $R_i$), controlled by $V_o$ (drawn as $+V_o-$ connected across the output terminals). Orange dashed boxes mark the ``A circuit'' ($R_i$, $A v_i$, $R_o$) and the ``$\beta$ circuit'' ($\beta V_o$ and $V_o$); orange arrows mark $R_{if}$ (looking into the loop from terminals $s$--$s'$) and $R_{of}$ (looking in from the output).}

\sekilnotu[7]{Two reduced circuits. Left: source $V_s$ ($+$ top, $-$ bottom; terminals $s$ and $s'$) with $R_s$ (struck out in red) in series with $R_{if}$. Right: dependent source $A_f V_s$ in series with $R_{of}$ (top wire, struck out in red); output $+v_o-$ across a resistor labelled $R_L$ (struck out in red) with the red label $R_{of}$ next to it.}

\begin{align*}
  A_f &= \frac{A}{1+A\beta} \\
  R_{if} &= (1+A\beta)\,R_i \\
  R_{of} &= \frac{R_o}{1+A\beta}
\end{align*}

\sayfa{8}
\section{2- Current Amplifiers (Series (o) -- Shunt (I))}
\[ A = \frac{I_o}{I_i} \]
\sekilnotu[8]{Block diagram of the series-shunt current amplifier: current source $I_s$ with parallel $R_s$ feeds the ``Basic current Amplifier'' block; the output current flows through the load $R_L$. The feedback network (``FBN'') takes the output current $I_o$ in series with the load path (orange oval and orange note ``Series'') and returns its output to the input node in parallel with $I_s$/$R_s$ (orange oval and orange note ``Shunt'').}

\sekilnotu[8]{Two-port equivalent: $I_s$ with $R_s$ and $R_{11}$ in parallel (arrow labelled $R_i$ looking into the input) feeds the ``Basic Circuit Amplifier'' box; at the output $R_{22}$ is in series with the load $R_L$, current $I_o$ (arrow) in the output loop, and an arrow labelled $R_o$ looking in at the output terminals.}

\sekilnotu[8]{Two sketches of the ``FBN'' box. Left: arrow labelled $R_{11}$ at the input side, output terminals open. Right: input terminals shorted together, arrow labelled $R_{22}$ at the output side.}

\sekilnotu[8]{``FBN'' box with input current $I_f$ entering and voltage $v_1$ ($+$, $-$) across the input terminals; a current source $I_o$ is connected at the output.}

\[ \beta = \frac{I_f}{I_o}\Big|_{v_1=0} \]

\sekilnotu[8]{Complete model: current source $I_s$ (terminals $s$ and $s'$) feeds $R_i$ (voltage $v_i$, $+$ top, $-$ bottom; input current $I_i$). Output side: dependent current source $A I_i$ in parallel with $R_o$, ending at the output terminals. Feedback part: dependent current source $\beta I_o$ connected across the input nodes (parallel with $R_i$), with the output current $I_o$ (arrow) flowing in a series path at the output terminals. Orange dashed boxes: ``A circuit'' ($R_i$, $A I_i$, $R_o$) and ``$\beta$ circuit'' ($\beta I_o$ and $I_o$); orange arrows mark $R_{if}$ (looking into $s$--$s'$) and $R_{of}$ (looking in from the output).}

\sayfa{9}
\sekilnotu[9]{Two small circuits: left, terminals $s$ and $s'$ with the current $I_s$ (arrow) entering $R_{if}$; right, dependent current source $A_f I_s$ in parallel with $R_{of}$, with the output current $I_o$ (arrow).}

\begin{align*}
  A_f &= \frac{A}{1+A\beta} \\
  R_{if} &= \frac{R_i}{1+A\beta} \\
  R_{of} &= R_o\,(1+A\beta)
\end{align*}

\[ \textcolor{turuncukalem}{1+\frac{R_2}{R_1} = 1001} \]

\section{Soru}
\sekilnotu[9]{Circuit: source $V_s$ with series $R_s$ feeds $R_{id}$ (voltage $v_1$, $+$ top, $-$ bottom; red arrow $R_{in}$ looking into it). Output side: dependent source $A_o v_1$ in series with $r_o$ to the output node $v_o$, with $R_L$ to ground. $R_1$ goes from the node below $R_{id}$ to ground, and $R_2$ connects this node to the output node. Annotations: large orange triangle (op-amp outline with $+$ and $-$) with the note ``Non inverting''; dashed orange ovals ``Shunt'' (around the output node) and ``series'' (around the $R_1$ node at the input); orange arrow $R_{of}$ and red arrow $R_{out}$ pointing at the output.}

\[ R_{in} = R_{if} - R_s \qquad\qquad R_{out} = \frac{1}{1/R_{of} - 1/R_L} \]

Find expressions for\\
$A$, $\beta$, $A_f$, $R_i$, $R_{if}$, $R_{of}$, and $R_o$\\
($A_o = 10^4$, $R_{id} = \qty{100}{\kilo\ohm}$, $r_o = \qty{1}{\kilo\ohm}$, $R_L = \qty{2}{\kilo\ohm}$, $R_1 = \qty{1}{\kilo\ohm}$, $R_2 = \qty{1}{\mega\ohm}$ and $R_s = \qty{10}{\kilo\ohm}$)\\
\underline{Feedback-type} : Shunt -- Series

\sekilnotu[9]{Two-port equivalent of the circuit. Input side: terminals with $V_s$ ($+$, $-$), $R_s$ on top, $R_{id}$ with $v_1$ ($+$ top, $-$ bottom) and a resistor $R_{11}$ in the bottom wire, labelled ``$R_{11}=??$'' and, in orange, $(R_1 \| R_2)$; arrow $R_i$ looking into the input. Output side: dependent source $A_o v_1$ in series with $r_o$, then $R_{22}$ and $R_L$ in parallel; orange note $(R_1+R_2)$ with an arrow pointing to $R_{22}$; arrow $R_o$ looking in at the output terminals.}

\sekilnotu[9]{Two small circuits of the feedback network ($R_1$, $R_2$). Left: $R_1$ from the top input node to the bottom wire, $R_2$ from the top node to a far end that is grounded (output shorted); arrow $R_{11}$ looking into the input. Right: input terminals open with $R_1$ from the top node to ground, $R_2$ from that node to the right output terminal; arrow $R_{22}$ looking into the output.}

\[ R_{11} = R_1 \parallel R_2 \qquad\qquad R_{22} = R_1 + R_2 \]

\sayfa{10}
\sekilnotu[10]{Circuit: $R_1$ in the left vertical branch with $v_f$ ($+$ top, $-$ bottom) across it; in the top wire a resistor (labelled $R_1$ in the note, presumably $R_2$) in series with a voltage source $v_o$ on the right.}

\[ \frac{v_f}{v_o} = \frac{R_1}{R_1+R_2} = \beta \]

\begin{align*}
  R_i &= R_s + R_{id} + R_{11} \\
  R_o &= r_o \parallel R_{22} \parallel R_L
\end{align*}

\[ {\color{kirmizikalem}A =}\ \underbrace{\frac{R_{id}}{R_i}}_{\displaystyle \frac{v_o}{v_s} = \frac{v_1}{v_s}} \cdot \underbrace{\frac{R_{22}\parallel R_L}{r_o + R_{22}\parallel R_L}}_{\displaystyle \frac{v_o}{v_1}\ \text{katkıları}} \cdot A_o \]

\[ A_f = \frac{A}{1+A\beta} \]

\[ A = \frac{\qty{100}{\kilo\ohm}}{\qty{111}{\kilo\ohm}} \cdot \frac{\qty{2}{\kilo\ohm}}{\qty{3}{\kilo\ohm}} \cdot 10^4 \approx 6000\ [\mathrm{V/V}] \]

\noindent
\begin{minipage}[t]{0.62\linewidth}
  \begin{align*}
    \beta &= \frac{R_1}{R_1+R_2} = 10^{-3}\ [\mathrm{V/V}] \\
    1 + A\beta &= 1 + 6000\cdot 10^{-3} = 7 \\
    A_f &= \frac{6000}{7} = 857\ [\mathrm{V/V}]
  \end{align*}
\end{minipage}\hfill
\begin{minipage}[t]{0.3\linewidth}
  $R_i = 111\mathrm{k}$\\[1mm]
  $R_o = 0{,}60\mathrm{k}$
\end{minipage}

\begin{align*}
  R_{if} &= (1+A\beta)\,\overset{111\mathrm{k}}{R_i} = \qty{777}{\kilo\ohm} \\
  R_{of} &= \frac{R_o}{(1+A\beta)} \approx \qty{93}{\ohm}
\end{align*}

\end{document}
